\documentclass[11pt,letterpaper]{article}

\usepackage[spanish,es-noshorthands]{babel}
\usepackage{fontspec}
\setmainfont{Source Serif Pro}
\usepackage{microtype}
\usepackage{geometry}
\usepackage{booktabs}
\usepackage{tabularx}
\usepackage{enumitem}
\usepackage{xcolor}
\usepackage{float}
\usepackage{csquotes}
\usepackage{amsmath}
\usepackage{siunitx}
\usepackage{tikz}
\usetikzlibrary{arrows.meta,positioning,fit,backgrounds}
\usepackage{pgfplots}
\pgfplotsset{compat=1.18}
\usepackage{xurl}
\usepackage[hidelinks]{hyperref}
\usepackage[backend=biber,style=apa,sorting=nyt,maxcitenames=2]{biblatex}
\DeclareLanguageMapping{spanish}{spanish-apa}

\geometry{margin=2.3cm}
% Los rótulos de las figuras usan nodos de ancho fijo y la bibliografía compone
% en modo sloppy: ambos generan avisos de línea holgada que son cosméticos.
\hbadness=10000
\setlength{\emergencystretch}{1em}
\setlength{\parindent}{0pt}
\setlength{\parskip}{0.42em}
\setlist{nosep}
\renewcommand{\topfraction}{0.92}
\renewcommand{\bottomfraction}{0.85}
\renewcommand{\textfraction}{0.06}
\renewcommand{\floatpagefraction}{0.72}
\setcounter{topnumber}{2}

\addbibresource{referencias.bib}
\AtBeginBibliography{\sloppy\small}
\setlength{\bibitemsep}{0.28em}
\setlength{\bibhang}{1.4em}

\sisetup{
  output-decimal-marker = {,},
  group-separator = {.},
  group-minimum-digits = 4,
  detect-all
}

% Colores: el granate del coste de secuenciar (el mismo tono que el gráfico
% del módulo), el azul de la ley de Moore, el verde de los datos oficiales del
% NHGRI, el ámbar de la producción mundial y el gris del texto secundario.
\definecolor{cCoste}{RGB}{163,31,52}
\definecolor{cMoore}{RGB}{42,120,190}
\definecolor{cNhgri}{RGB}{36,122,82}
\definecolor{cProd}{RGB}{176,110,20}
\definecolor{cEje}{RGB}{26,58,102}
\definecolor{cGris}{RGB}{120,116,110}

\pgfplotsset{
  informe/.style={
    width=\textwidth, height=7.2cm,
    tick label style={font=\footnotesize},
    label style={font=\footnotesize},
    legend style={font=\footnotesize, draw=cGris!50, fill=white,
      fill opacity=0.92, text opacity=1, cells={anchor=west}},
    grid=major, major grid style={cGris!22},
    axis line style={cGris!80},
    /pgf/number format/.cd, use comma, 1000 sep={.},
  },
}

\tikzset{
  every node/.append style={execute at begin node={\hyphenpenalty=10000}},
  % Celda de la matriz 5x5: #1 es el color de relleno
  celda/.style={draw=cGris!60, fill=#1, minimum width=2.85cm,
    minimum height=1.32cm, text width=2.6cm, align=center,
    font=\scriptsize, inner sep=2pt},
  cabecera/.style={font=\footnotesize\bfseries, text=cEje, align=center},
}

% Celda de la matriz: #1 columna (1 a 5), #2 fila (1 a 5), #3 estilo de
% relleno, #4 texto
\newcommand{\celda}[4]{%
  \node[celda=#3] (c#1#2) at ({2.4+(#1-0.5)*2.85},{-(#2-0.5)*1.32}) {#4};}

\begin{document}

% --- Portada -----------------------------------------------------------------
\begin{titlepage}
\centering
\vspace*{3.2cm}
{\large\scshape Gestión del Desarrollo Tecnológico:\\
Estrategias y Análisis de Portfolio\par}
\vspace{0.5cm}
{\normalsize Módulo 5. Actividad 5.3: Posición del ADN en la matriz tecnológica\par}
\vspace{2.6cm}
{\color{cEje}\rule{0.55\textwidth}{0.8pt}\par}
\vspace{0.7cm}
{\LARGE\bfseries Leer la vida a precio de saldo\par}
\vspace{0.5cm}
{\large Dónde cae la secuenciación del ADN en la matriz 5\,×\,5\\
y a qué ritmo se ha abaratado desde el año 2000\par}
\vspace{0.7cm}
{\color{cEje}\rule{0.55\textwidth}{0.8pt}\par}
\vfill
{\large Carlos Luis Rojas Aragonés\par}
\vspace{0.3cm}
{\normalsize Chief Technology Officer Program\par}
\vspace{0.3cm}
{\normalsize 5 de octubre de 2026\par}
\vspace*{1.5cm}
\end{titlepage}

% --- Cuerpo ------------------------------------------------------------------
\section{Las dos respuestas, por adelantado}
\label{sec:respuestas}

\textbf{Pregunta 1. ¿En qué celda de la matriz 5\,×\,5 coloco la secuenciación
del ADN?} En la intersección del operando \textbf{información} con el proceso
\textbf{transformar}. Un secuenciador toma una molécula física que pertenece a
un organismo vivo y la convierte en un dato digital: una cadena de letras A, C,
G y T con su calidad asociada. El organismo no sale modificado; lo que se crea
es información. La celda vecina de \emph{organismos vivos por transformar} le
corresponde a la tecnología complementaria, la edición genética, y la
secuenciación queda justo en la frontera entre ambas columnas
(sección~\ref{sec:matriz}).

\textbf{Pregunta 2. ¿Cuál es la tasa anual de mejora del coste de secuenciar un
genoma humano completo desde el año 2000?} Leída en el gráfico del módulo, la
curva de coste baja de unos \num{2,2e5} en el año 2000 a algo más de 1 en
2019, en la misma escala: un factor de \num{192000} en diecinueve años. Eso es
una \textbf{mejora de aproximadamente el 90~\% anual}. Dicho de las otras dos
maneras equivalentes, \textbf{el coste se multiplica cada año por 0,53} (cae un
47~\%) y \textbf{se reduce a la mitad cada 13 meses}. Es más del doble del
ritmo de la ley de Moore \parencite{moore1965}, que en el mismo gráfico mejora un 40~\% anual y se
reduce a la mitad cada dos años (sección~\ref{sec:tasa}).

Los datos oficiales del NHGRI, que no salen del gráfico sino de los costes
reales de sus centros de secuenciación, confirman la cifra de manera
independiente: un 93~\% anual entre 2001 y 2019 y un 80~\% anual si se extiende
hasta 2022, cuando el genoma costaba 525 dólares frente a los 95 millones de
2001 (sección~\ref{sec:nhgri}). La media, sin embargo, esconde tres épocas muy
distintas, y esa es la conclusión que más le importa a quien gestiona una
cartera tecnológica (sección~\ref{sec:eras}).

\section{\texorpdfstring{La matriz tecnológica 5\,×\,5}{La matriz tecnológica 5 × 5}}
\label{sec:matriz}

\subsection*{Qué clasifica la matriz}

La matriz clasifica las tecnologías por su función, no por su disciplina ni por
su industria. \textcite{deweck2022} parte de una idea simple: toda tecnología
ejecuta al menos un \emph{proceso} sobre al menos un \emph{operando}, el objeto
que cambia de estado. La versión clásica de \textcite{vanwyk1988} era una
cuadrícula de 3\,×\,3, con materia, energía e información como operandos y
transformar, transportar y almacenar como procesos. De Weck la amplía a
5\,×\,5 añadiendo dos operandos, los \textbf{organismos vivos} y el
\textbf{valor} (el dinero y los activos financieros), y dos procesos,
\textbf{intercambiar} y \textbf{controlar}. La figura~\ref{fig:matriz} la
reproduce con un ejemplo por celda y marca dónde cae la secuenciación y las
tecnologías genómicas que la rodean.

\begin{figure}[htbp]
\centering
\begin{tikzpicture}
  % cabeceras de columna (operandos) y de fila (procesos)
  \foreach \t [count=\j] in {Materia, Energía, Información, Organismos\\vivos, Valor}
    \node[cabecera] at ({2.4+(\j-0.5)*2.85},0.55) {\t};
  \foreach \t [count=\i] in {Transformar, Transportar, Almacenar, Intercambiar, Controlar}
    \node[cabecera, anchor=east] at (2.25,{-(\i-0.5)*1.32}) {\t};
  \node[font=\scriptsize\itshape, text=cGris, anchor=south] at (9.525,1.55)
    {operando: el objeto que cambia de estado};
  \node[font=\scriptsize\itshape, text=cGris, rotate=90, anchor=south]
    at (0.15,-3.3) {proceso};

  % fila 1: transformar
  \celda{1}{1}{cCoste!14}{Síntesis de ADN,\\fundición, impresión 3D}
  \celda{2}{1}{white}{Motor, turbina,\\panel solar}
  \celda{3}{1}{cCoste!80}{\color{white}\bfseries Secuenciación\\de ADN}
  \celda{4}{1}{cCoste!14}{Edición genética\\(CRISPR), cirugía}
  \celda{5}{1}{white}{Banca de inversión,\\emisión de moneda}
  % fila 2: transportar
  \celda{1}{2}{white}{Camión, tubería,\\cinta transportadora}
  \celda{2}{2}{white}{Red eléctrica\\de transmisión}
  \celda{3}{2}{cCoste!14}{Fibra óptica, nube\\(datos genómicos)}
  \celda{4}{2}{white}{Ambulancia,\\avión de pasajeros}
  \celda{5}{2}{white}{Transferencia\\interbancaria}
  % fila 3: almacenar
  \celda{1}{3}{white}{Silo, almacén,\\contenedor}
  \celda{2}{3}{white}{Batería, embalse}
  \celda{3}{3}{cCoste!14}{Disco, archivo de\\genomas, datos en ADN}
  \celda{4}{3}{cCoste!14}{Biobanco,\\criopreservación}
  \celda{5}{3}{white}{Banco, cartera\\de activos}
  % fila 4: intercambiar
  \celda{1}{4}{white}{Puerto, mercado\\de materias primas}
  \celda{2}{4}{white}{Mercado eléctrico\\mayorista}
  \celda{3}{4}{cCoste!14}{GenBank, GISAID,\\Internet}
  \celda{4}{4}{white}{Trasplante,\\banco de semillas}
  \celda{5}{4}{white}{Bolsa de valores}
  % fila 5: controlar
  \celda{1}{5}{white}{Válvula, robot\\industrial}
  \celda{2}{5}{white}{Regulador de red,\\termostato}
  \celda{3}{5}{white}{Software de control,\\criptografía}
  \celda{4}{5}{cCoste!14}{Medicina de precisión,\\marcapasos}
  \celda{5}{5}{white}{Banco central,\\auditoría}

  % la celda elegida y la flecha que muestra el cambio de operando
  \draw[cCoste, line width=1.6pt] (c31.south west) rectangle (c31.north east);
  \draw[-{Stealth[length=2.4mm]}, line width=1.1pt, cEje]
    (11.25,1.0) to[bend right=28]
    node[midway, above, font=\scriptsize\bfseries, text=cEje] {lee}
    (8.95,1.0);
\end{tikzpicture}
\caption{Matriz tecnológica 5\,×\,5 de operandos por procesos, según la
estructura de \textcite{deweck2022}, con un ejemplo propio en cada celda. La
celda granate es la que propongo para la secuenciación del ADN; la flecha
indica que su materia prima procede de la columna de organismos vivos y su
producto cae en la de información. Las celdas rosadas son las tecnologías
genómicas que la rodean.}
\label{fig:matriz}
\end{figure}

\subsection*{Por qué información por transformar, y no otra celda}

La regla para ubicar una tecnología es preguntarse qué objeto sale de ella con
un estado distinto del que tenía al entrar. Aplicada al secuenciador, la
respuesta es inequívoca: entra una muestra de ADN y sale un archivo. El
cuadro~\ref{tab:celdas} recorre las cuatro celdas que podrían reclamar la
secuenciación y explica por qué tres de ellas no encajan.

\begin{table}[htbp]
\centering
\small
\renewcommand{\arraystretch}{1.25}
\begin{tabularx}{\textwidth}{@{}p{3.6cm}X p{1.9cm}@{}}
\toprule
\textbf{Celda candidata} & \textbf{Argumento} & \textbf{Veredicto} \\
\midrule
Información por transformar &
  El producto es un dato nuevo, la secuencia, que antes no existía en forma
  legible. Sus figuras de mérito son de información: coste por genoma, bases
  leídas por día, error por base (puntuación Phred) y longitud de lectura. &
  \textbf{Celda principal} \\
Organismos vivos por transformar &
  La muestra viene de un organismo, pero el organismo no cambia. Esta celda
  pertenece a la edición genética, que sí altera al ser vivo. &
  Frontera, no celda \\
Información por almacenar &
  El ADN es el soporte de información más denso que se conoce y existe
  investigación para guardar datos en él, pero esa es otra tecnología. El
  secuenciador lee el soporte; no lo usa para guardar. &
  Celda vecina \\
Materia por transformar &
  Hay química en el proceso (amplificación, fluoróforos, nanoporos), pero es
  el medio, no el fin. Nadie compra un secuenciador por los reactivos que
  consume. &
  Descartada \\
\bottomrule
\end{tabularx}
\caption{Celdas candidatas para la secuenciación del ADN.}
\label{tab:celdas}
\end{table}

La matriz también revela algo útil para la gestión: la secuenciación es la
\emph{puerta de entrada} de un ecosistema que ocupa media columna de
información. Cada genoma secuenciado hay que transportarlo, almacenarlo e
intercambiarlo (filas 2 a 4 de esa columna), y sólo entonces se convierte en
control sobre un organismo vivo, que es lo que hace la medicina de precisión.
Cuando el coste de la primera celda se desploma, como muestra la sección
siguiente, el cuello de botella se desplaza a las demás. Un genoma humano a
30 veces de cobertura ocupa del orden de un centenar de gigabytes de datos
crudos; con secuenciar barato, almacenar y analizar se vuelve la parte cara.

\section{La tasa anual de mejora desde el año 2000}
\label{sec:tasa}

\subsection*{Método}

Una figura de mérito que progresa de forma exponencial cumple
$\mathrm{FOM}(t) = \mathrm{FOM}_0\,(1+r)^{t-t_0}$, donde $r$ es la tasa anual
de mejora \parencite{deweck2022}. El coste es una figura de mérito en la que
menos es mejor, así que se trabaja con su inverso, los genomas que se pueden
secuenciar por dólar. Con los costes $C_0$ en el año $t_0$ y $C_1$ en el año
$t_1$:
\begin{equation}
  r = \left(\frac{C_0}{C_1}\right)^{1/(t_1-t_0)} - 1,
  \qquad
  d = 1 - \left(\frac{C_1}{C_0}\right)^{1/(t_1-t_0)},
  \qquad
  T_{1/2} = \frac{\ln 2}{\ln(1+r)}.
  \label{eq:tasa}
\end{equation}
$r$ es la mejora anual, $d$ la reducción anual del coste y $T_{1/2}$ el tiempo
que tarda el coste en reducirse a la mitad. Las tres expresan el mismo dato.

Hay una ventaja importante en este método. El gráfico del módulo es una
adaptación del de \textcite{shendure2017} y su rotulación es ambigua: el título
central dice \enquote{coste por genoma humano} y el eje derecho, en el mismo
color, dice \enquote{coste por MB}. Los valores absolutos tampoco coinciden con
ninguna de las dos magnitudes en las series oficiales (en 2012 la curva marca
unos 2 dólares). Para la tasa da igual: la ecuación~\eqref{eq:tasa} sólo usa
el \emph{cociente} entre dos lecturas, y un cambio de escala o de unidades
multiplica numerador y denominador por el mismo factor. La pendiente en escala
logarítmica es la misma se lea el eje que se lea.

\subsection*{Lectura del gráfico}

Para no depender de una lectura a ojo, digitalicé las tres curvas del gráfico
píxel a píxel, calibrando los ejes con las marcas de graduación (70,2 píxeles
por década en vertical y 52,4 por año en horizontal). La
figura~\ref{fig:grafico} reconstruye la curva de coste con esos datos y la
compara con la línea de la ley de Moore del propio gráfico.

\begin{figure}[htbp]
\centering
\begin{tikzpicture}
\begin{axis}[informe,
  ymode=log, xmin=1995, xmax=2019.2, ymin=0.4, ymax=4e9,
  xtick={1996,1998,...,2018},
  /pgf/number format/1000 sep={},
  ylabel={Coste en la escala del gráfico (dólares)},
  legend pos=north east,
]
\addplot[cMoore, thick, dashed, domain=1995:2019, samples=2]
  {1e9*10^(-(x-1995)*3.51/24)};
\addlegendentry{Ley de Moore: $\times 0{,}71$ por año}
\addplot[cCoste, very thick] table[x=anio, y=valor] {datos/grafico-coste.dat};
\addlegendentry{Coste de secuenciar (digitalizado)}
\addplot[black, thick, densely dotted, mark=*, mark size=2.2pt,
  mark options={solid, fill=white}]
  coordinates {(2000,218872) (2019,1.14)};
\addlegendentry{Tasa media 2000-2019: $\times 0{,}53$ por año}
\node[font=\scriptsize, anchor=south west, align=left] at (axis cs:2000.2,3e5)
  {2000: \num{2,19e5}};
\node[font=\scriptsize, anchor=south east, align=right] at (axis cs:2018.8,1.6)
  {2019: 1,14};
\draw[cEje, -{Stealth}] (axis cs:2004.6,1.2e2) -- (axis cs:2007.8,1.2e2);
\node[font=\scriptsize, text=cEje, align=left, anchor=east]
  at (axis cs:2004.5,1.2e2) {454 e Illumina:\\caída de 8.000 veces\\en seis años};
\end{axis}
\end{tikzpicture}
\caption{Curva de coste del gráfico del módulo, digitalizada, frente a la ley
de Moore. La recta punteada une las lecturas de 2000 y 2019; su pendiente es la
tasa media que pide la actividad. Fuente: elaboración propia a partir del
gráfico del módulo, adaptado de \textcite{shendure2017}.}
\label{fig:grafico}
\end{figure}

El cuadro~\ref{tab:calculo} aplica la ecuación~\eqref{eq:tasa} a las lecturas.
La tasa media de 2000 a 2019 es del \textbf{89,7~\% anual}. Para acotar el
error de lectura supuse un margen de una décima de década (un 26~\% del valor)
en cada extremo y en el sentido desfavorable: la tasa queda entre el
\textbf{85~\% y el 94~\%}. Un ajuste por mínimos cuadrados en escala
logarítmica sobre todos los puntos de 2000 a 2019 da una cifra más alta, del
111~\%, porque la caída abrupta de la mitad del periodo pesa más que los
extremos planos; prefiero la tasa entre extremos porque es la que responde a
la pregunta, que es cuánto mejoró el coste \emph{desde} el año 2000 hasta hoy.

\begin{table}[htbp]
\centering
\small
\renewcommand{\arraystretch}{1.2}
\begin{tabular}{@{}lrrrrr@{}}
\toprule
\textbf{Periodo} & \textbf{Factor} & \textbf{Años} & \textbf{Mejora $r$} &
  \textbf{Caída $d$} & \textbf{$T_{1/2}$} \\
\midrule
\textbf{2000-2019 (respuesta)} & \num{192000} & 19 & \textbf{89,7~\%} &
  47,3~\% & 13 meses \\
\quad Margen de lectura & & & 85-94~\% & 46-49~\% & 12-14 meses \\
\quad Ajuste logarítmico & & & 111~\% & 53~\% & 11 meses \\
\midrule
Ley de Moore (mismo gráfico) & \num{3200} & 24 & 40,0~\% & 28,6~\% &
  25 meses \\
\bottomrule
\end{tabular}
\caption{Tasa anual de mejora del coste, calculada con la
ecuación~\eqref{eq:tasa} sobre la curva digitalizada.}
\label{tab:calculo}
\end{table}

Una precisión sobre \enquote{genoma humano completo}. El genoma que se
publicó en 2003 todavía tenía alrededor de un 8~\% sin resolver, sobre todo
centrómeros y regiones repetitivas, y la primera secuencia de verdad completa,
de telómero a telómero, no llegó hasta 2022 \parencite{nurk2022}. El coste que
mide el gráfico, y el que mide el NHGRI, es el de un genoma con calidad de
referencia según los estándares de cada época, que es la convención habitual.

\section{Contraste con los datos oficiales del NHGRI}
\label{sec:nhgri}

El National Human Genome Research Institute publica desde 2001 el coste real
de secuenciar un genoma humano en los centros que financia, con todos los
gastos incluidos: reactivos, personal, amortización de equipos, gestión de
datos y gastos indirectos \parencite{nhgri2023}. Es la fuente canónica de esta
curva, y como no procede del gráfico del módulo sirve como comprobación
independiente (figura~\ref{fig:nhgri}).

\begin{figure}[htbp]
\centering
\begin{tikzpicture}
\begin{axis}[informe,
  ymode=log, xmin=2001, xmax=2023.3, ymin=40, ymax=3e8,
  xtick={2001,2003,...,2023},
  /pgf/number format/1000 sep={},
  ylabel={Coste por genoma humano (dólares)},
  legend pos=north east,
]
\addplot[cMoore, thick, dashed, domain=2001.708:2023, samples=2]
  {95263071.9*0.5^((x-2001.708)/2)};
\addlegendentry{Si hubiera seguido la ley de Moore}
\addplot[cNhgri, very thick, mark=*, mark size=1.1pt]
  table[x=anio, y=genoma] {datos/nhgri-genoma.dat};
\addlegendentry{Coste real, NHGRI}
\addplot[only marks, mark=square*, mark size=2.2pt, cCoste]
  coordinates {(2022.4,100) (2022.75,200)};
\addlegendentry{Costes anunciados por fabricantes}
\node[font=\scriptsize, anchor=west, text=cNhgri] at (axis cs:2001.9,9.5e7)
  {2001: 95,3 millones};
\node[font=\scriptsize, anchor=south, text=cNhgri, align=center]
  at (axis cs:2021.2,1.4e3) {2022: 525 \$};
\node[font=\scriptsize, anchor=east, text=cCoste, align=center]
  at (axis cs:2022.1,140) {200 y 100 \$};
\draw[cEje, -{Stealth}] (axis cs:2010.1,2e5) -- (axis cs:2008.5,5e5);
\node[font=\scriptsize, text=cEje, anchor=west, align=left]
  at (axis cs:2010.2,2e5) {2008: llega la secuenciación\\de nueva generación};
\end{axis}
\end{tikzpicture}
\caption{Coste por genoma humano según el NHGRI (77 observaciones de 2001 a
2022) frente a una trayectoria hipotética que se redujera a la mitad cada dos
años desde el mismo punto de partida. Los cuadrados son los costes anunciados
en 2022 por \textcite{illumina2022} y \textcite{ultima2022}, que no son
costes totales comparables con la serie. Fuentes: \textcite{nhgri2023};
\textcite{datahub2023}.}
\label{fig:nhgri}
\end{figure}

Los resultados coinciden con el gráfico del módulo. Entre septiembre de 2001 y
principios de 2019 la mejora es del \textbf{93,2~\% anual}, a tres puntos y
medio de la lectura del gráfico. Si se extiende hasta mayo de 2022, cuando la
serie se detiene en 525 dólares, la tasa baja al \textbf{79,7~\% anual},
porque se añaden tres años de la fase más lenta. Ambas cifras confirman el
orden de magnitud de la respuesta: una mejora del 80 al 90~\% anual, el doble
largo de la ley de Moore. La recta azul de la figura~\ref{fig:nhgri} pone ese
contraste en dinero: con el ritmo de Moore, un genoma costaría hoy unos
\num{74000} dólares, no 525.

\section{Una media que esconde tres épocas}
\label{sec:eras}

La tasa media es la respuesta correcta a la pregunta, pero no es una buena
herramienta de planificación, porque el coste no bajó a ritmo constante. Bajó
en tres escalones, que se ven en las dos curvas y que el
cuadro~\ref{tab:eras} y la figura~\ref{fig:eras} separan.

\begin{table}[htbp]
\centering
\small
\renewcommand{\arraystretch}{1.2}
\begin{tabularx}{\textwidth}{@{}p{3.2cm}Xrr@{}}
\toprule
\textbf{Época} & \textbf{Tecnología dominante} & \textbf{Gráfico} &
  \textbf{NHGRI} \\
\midrule
Sanger capilar & Electroforesis capilar automatizada (ABI 3730): un fragmento
  por capilar, 96 a la vez. Mejora incremental de un diseño dominante. &
  76~\% & 53~\% \\
Nueva generación & 454, Illumina GA, GAII y HiSeq: millones de fragmentos en
  paralelo sobre una sola superficie. Cambio de arquitectura. &
  348~\% & 451~\% \\
Maduración & Química de Illumina consolidada; mejoras de rendimiento por
  equipo y competencia de nuevos fabricantes al final del periodo. &
  10~\% & 29~\% \\
\midrule
\textbf{Periodo completo} & & \textbf{90~\%} & \textbf{80~\%} \\
\bottomrule
\end{tabularx}
\caption{Mejora anual del coste en cada época. Cortes del gráfico: 2000-2004,
2004-2010 y 2010-2019. Cortes del NHGRI, con su propia cronología:
2001-2007, 2007-2011 y 2011-2022.}
\label{tab:eras}
\end{table}

\begin{figure}[htbp]
\centering
\begin{tikzpicture}
\begin{axis}[informe, height=6.6cm,
  ybar, bar width=17pt, ymin=0, ymax=520,
  symbolic x coords={Sanger, NGS, Maduración, Total},
  xticklabels={Sanger capilar, Nueva generación, Maduración, Periodo completo},
  xtick=data, enlarge x limits=0.18,
  ylabel={Mejora anual del coste (\%)},
  nodes near coords, every node near coord/.append style={font=\scriptsize,
    fill=white, inner sep=1pt},
  legend style={at={(0.98,0.97)}, anchor=north east},
  xmajorgrids=false,
]
\addplot[fill=cCoste, draw=cCoste] coordinates
  {(Sanger,76.4) (NGS,347.6) (Maduración,10.5) (Total,89.7)};
\addlegendentry{Gráfico del módulo}
\addplot[fill=cNhgri, draw=cNhgri] coordinates
  {(Sanger,53.1) (NGS,451.2) (Maduración,29.0) (Total,79.7)};
\addlegendentry{NHGRI}
\draw[draw=cMoore, line width=0.8pt, dash pattern=on 3pt off 2pt] ({rel axis cs:0,0}|-{axis cs:Sanger,40})
  -- ({rel axis cs:1,0}|-{axis cs:Sanger,40});
\addlegendimage{line legend, cMoore, thick, dashed}
\addlegendentry{Ley de Moore (40~\%)}
\end{axis}
\end{tikzpicture}
\caption{Mejora anual del coste por época. Sólo la maduración posterior a 2010
queda por debajo de la ley de Moore; la época de nueva generación la supera
por un factor de diez.}
\label{fig:eras}
\end{figure}

Las tres épocas son las tres fases de una sucesión de curvas S. La primera es
la mejora incremental de una tecnología dominante, la electroforesis capilar,
que ya mejoraba más deprisa que los chips. La segunda es un \emph{salto de
curva}: la secuenciación masiva en paralelo no hizo más rápido lo mismo, sino
que cambió la arquitectura, y en seis años redujo el coste unas \num{8000}
veces según el gráfico. La tercera es la parte alta de la nueva curva S, con
rendimientos decrecientes de la química de Illumina, hasta que los anuncios de
2022 de 200 y 100 dólares por genoma señalan el inicio de otro posible salto
\parencite{illumina2022,ultima2022}.

Para un CTO esto tiene una lectura práctica. Si en 2010 alguien hubiera
extrapolado la tasa de la época anterior para presupuestar un programa
genómico a diez años, habría sobrestimado el coste final de forma enorme. Si
hubiera extrapolado la tasa de 2004 a 2010, lo habría subestimado por varios
órdenes de magnitud. La tasa media de 90~\% anual es correcta como historia,
pero la decisión depende de saber en qué tramo de la curva S está la
tecnología hoy.

\section{Material de apoyo: la otra cara de la misma curva}
\label{sec:produccion}

El gráfico del módulo tiene una tercera curva que la pregunta no usa: la
producción mundial de secuenciación, en kilobases por día. La digitalicé con
el mismo procedimiento. Pasó de unos 56 kilobases diarias en el año 2000 a
unos 440 millones en 2014: una mejora del \textbf{211~\% anual}, que duplica la
capacidad mundial cada siete meses, tres veces más rápido de lo que cayó el
coste.

Las dos curvas están ligadas. La figura~\ref{fig:wright} enfrenta el coste con
la producción en los años en que el gráfico tiene ambas, de 2000 a 2014, en
escala logarítmica en los dos ejes. Los puntos caen casi sobre una recta
($R^2 = 0{,}995$) de pendiente $-0{,}77$: cada vez que la producción mundial
se multiplica por diez, el coste se divide por 5,9. Es la forma de la ley de
\textcite{wright1936}, la curva de experiencia, aunque con una advertencia
metodológica: Wright usa la producción \emph{acumulada}, y aquí sólo tengo el
ritmo diario, así que la relación describe el acoplamiento entre volumen y
coste pero no prueba en qué sentido va la causa. Lo más probable es que vaya
en los dos: el abaratamiento abrió usos nuevos y el volumen financió la
siguiente generación de equipos, el ciclo que \textcite{carlson2003} describió
para las tecnologías biológicas cuando la segunda época apenas empezaba.

\begin{figure}[htbp]
\centering
\begin{tikzpicture}
\begin{axis}[informe, height=6.8cm,
  xmode=log, ymode=log, xmin=20, xmax=2e9, ymin=0.5, ymax=1e6,
  xlabel={Producción mundial (kilobases por día)},
  ylabel={Coste en la escala del gráfico},
  legend pos=north east,
]
\addplot[only marks, mark=*, mark size=1.5pt, cProd]
  table[x=produccion, y=coste] {datos/coste-vs-produccion.dat};
\addlegendentry{Años 2000 a 2014}
\addplot[cCoste, thick, domain=40:1e9, samples=2] {10^(6.491)*x^(-0.7713)};
\addlegendentry{Ajuste: $C \propto P^{-0{,}77}$, $R^2=0{,}995$}
\node[font=\scriptsize, anchor=south west] at (axis cs:60,2.4e5) {2000};
\node[font=\scriptsize, anchor=west] at (axis cs:6e8,1.5) {2014};
\end{axis}
\end{tikzpicture}
\caption{Coste de secuenciar frente a producción mundial, ambos leídos del
gráfico del módulo. La relación es casi exactamente una ley de potencia.}
\label{fig:wright}
\end{figure}

\section{Conclusiones}
\label{sec:conclusiones}

\begin{enumerate}[leftmargin=1.4em, itemsep=0.4em]
  \item \textbf{La secuenciación del ADN pertenece a la celda de información
    por transformar}, en la frontera con la columna de organismos vivos: lee
    un objeto biológico y produce información. La edición genética es su
    complemento en la celda de al lado, y la columna de información que queda
    por debajo (transportar, almacenar, intercambiar) es donde se desplaza el
    cuello de botella cuando secuenciar se vuelve barato.
  \item \textbf{La tasa anual de mejora del coste desde el año 2000 es de
    aproximadamente el 90~\%} según el gráfico del módulo (entre 85 y 94~\%
    con el margen de lectura): el coste cae un 47~\% cada año y se reduce a la
    mitad cada 13 meses. Los datos del NHGRI la confirman con un 93~\% hasta
    2019 y un 80~\% hasta 2022.
  \item \textbf{Es más del doble del ritmo de la ley de Moore}, que en el mismo
    gráfico mejora un 40~\% anual. Con ese ritmo, un genoma costaría hoy unos
    \num{74000} dólares en lugar de 525.
  \item \textbf{La media resume tres épocas muy distintas}: 76~\% anual con
    Sanger, 350~\% con la secuenciación de nueva generación y apenas 10~\% en
    la maduración posterior a 2010. Para planificar no sirve la media, sirve
    saber en qué tramo de la curva S se está.
\end{enumerate}

\printbibliography[title={Referencias}]

\end{document}
